% Shared exercise chunk: l4-latent-likelihood
\begin{exercisechunk}
\item \textbf{Latent likelihood and the forward algorithm.}
\label[exercise]{ex:latent-likelihood}
\exerciseprofile{2}{2}{2}{\exproof\quad\excalculation\quad\exinterpretation}{%
Derive the forward recursion and explain why it replaces enumeration of
$d^{T+1}$ hidden trajectories by $O(Td^2)$ computation.}
Let \(d,N,T\) be positive integers. Let
\(\pi\in[0,1]^d\) be a probability vector, and let
\(A\in[0,1]^{d\times d}\) and \(C\in[0,1]^{N\times d}\) satisfy
\[
 \sum_{i=1}^d A_{ij}=1,
 \qquad
 \sum_{a=1}^N C_{aj}=1
 \quad\text{for every }j\in[d].
\]
Let \(Q\) be the joint probability mass function of a hidden Markov model on
\[
 [d]^{T+1}\times[N]^T.
\]
The initial-state pmf is \(\pi\). Given previous state \(j\), the next-state
pmf is the column \(A_{\mathord\cdot j}\). Given current state \(j\), the
emission pmf is the column \(C_{\mathord\cdot j}\). These conditional pmfs do
not depend on \(t\).
Fix \(x_{1:T}\in[N]^T\). For \(s_{0:T}\in[d]^{T+1}\), write
\(Q(s_{0:T},x_{1:T})\) for the joint mass and \(Q(x_{1:T})\) for the
marginal mass of \(x_{1:T}\).
\begin{enumerate}[(a),beginpenalty=10000]
\item Derive
\[
 Q(s_{0:T},x_{1:T})=\pi_{s_0}
 \prod_{t=1}^T A_{s_ts_{t-1}}C_{x_ts_t},
 \qquad
 Q(x_{1:T})=\sum_{s_{0:T}\in[d]^{T+1}}
 \pi_{s_0}\prod_{t=1}^T A_{s_ts_{t-1}}C_{x_ts_t}.
\]
Verify that the sum has \(d^{T+1}\) terms.
\item For \(t\in[T]\), let \(Q_t\) be the joint marginal pmf on
\([N]^t\times[d]\). Define \(\Lambda_0=\pi\) and
\(\Lambda_t\in\mathbb R^d\) by
\(\Lambda_t(j)=Q_t(x_{1:t},j)\) for \(j\in[d]\).
Derive a recursion for \(\Lambda_t\) and show that, in vector form, it is
\[
 \Lambda_t=\diag(C_{x_t,1},\ldots,C_{x_t,d})A\Lambda_{t-1}.
\]
\item Let \(\mathbf1\in\mathbb R^d\) be the all-ones vector. Show that
\(Q(x_{1:T})=\mathbf1^\top\Lambda_T\) and that the recursion
evaluates it in \(O(Td^2)\) arithmetic operations, rather than summing
\(d^{T+1}\) terms.
\end{enumerate}
\end{exercisechunk}
